%\title{Apache Maven} 
%\author{Niederegger Andreas }
%\date{Mai 2015}

%\maketitle

\section{Apache Maven}
\subsection{Einleitung}
Maven ist ein Build-Management-Tool der Apache Software Foundation und basiert auf Java. Mit ihm kann man insbesondere Java-Programme standardisiert erstellen und verwalten.
Der Name Maven kommt aus dem Jiddischen und bedeutet so viel wie "Sammler des Wissens“.


\begin{figure}[ht]
  \centering
    \includegraphics[width=0.4\textwidth]{MavenLogo}
  \caption{Maven Logo}
  \label{fig:MavenLogo}
\end{figure}


\subsection{Design}
Maven basiert auf einer Plugin-Architektur, die es ermöglicht, Plugins für verschiedene Anwendungen auf das Projekt anzuwenden, ohne diese explizit installieren zu müssen. 
\\Die Bandbreite reicht von Plugins, die es ermöglichen, direkt aus Maven heraus eine Webanwendung zu starten, um sie im Browser zu testen, über welche, die es ermöglichen, Datenbanken zu testen oder zu erstellen, bis hin zu solchen, die Web Services generieren. Die dafür nötigen Tätigkeiten beschränken sich häufig nur darauf, das gewünschte Plugin zu ermitteln und einzusetzen.
\newpage
\subsection{Lebenszyklus}
Der Standard-Lebenszyklus
Maven geht von einem Zyklus aus, der bei der Softwareerstellung häufig durchlaufen wird. Dabei wird nicht unterstellt, dass jedes Softwareprojekt alle Phasen des im Folgenden dargestellten Zyklus verwendet:

\begin{itemize}
\item archetype (Scaffolding): Damit kann ein Template für ein Softwareprojekt erstellt werden. Abhängigkeiten werden aufgelöst und bei Bedarf heruntergeladen.
\item validate (Validieren): Hier wird geprüft, ob die Projektstruktur gültig und vollständig ist.
\item compile (Kompilieren): In dieser Phase wird der Quellcode kompiliert.
\item test (Testen): Hier wird der kompilierte Code mit einem passenden Testframework getestet. Maven berücksichtigt dabei in späteren Zyklen, dass Testklassen normalerweise nicht in der auszuliefernden Software vorhanden sind.
\item package (Verpacken): Das Kompilat wird – ggf. mit anderen nichtkompilierbaren Dateien – zur Weitergabe verpackt. Häufig handelt es sich dabei um eine Jar-Datei.
\item integration-test (Test der Integrationsmöglichkeit): Das Softwarepaket wird auf eine Umgebung (anderer Rechner, anderes Verzeichnis, Anwendungsserver) geladen und seine Funktionsfähigkeit geprüft.
\item verify (Gültigkeitsprüfung des Software-Pakets): Prüfungen, ob das Softwarepaket eine gültige Struktur hat und ggf. bestimmte Qualitätskriterien erfüllt.
\item install (Installieren im lokalen Maven-Repository): Installiere das Softwarepaket in dem lokalen Maven-Repository, um es in anderen Projekten verwenden zu können, die von Maven verwaltet werden. Dies ist insbesondere für modulare Projekte interessant.
\item deploy (Installieren im fernen Maven-Repository): Stabile Versionen der Software werden auf einem fernen Maven-Repository installiert und stehen damit in Umgebungen mit mehreren Entwicklern allen zur Verfügung.
\end{itemize}


\subsection{Unterstützte Entwicklungsumgebungen}
Für die gängigsten Entwicklungsumgebungen (z. B. Eclipse, IntelliJ IDEA oder NetBeans) sind Plugins vorhanden, über die sich Maven direkt aus der Entwicklungsumgebung heraus bedienen lässt. Zusätzlich sind Maven-Plugins vorhanden, die Dateien erzeugen, welche den Import eines reinen Maven-Projekts in die Entwicklungsumgebung ermöglichen 

\subsection{Teilprojekte}
Die Entwicklung von Maven ist in verschiedene Teilprojekte untergliedert.
\begin{itemize}
\item Maven 1.x pflegt die älteren Versionen von Maven.
\item Maven 2 entwickelt die aktuelle Produktlinie von Maven weiter.
\item Plugins entwickelt die meisten Maven-Plugins.
\item Shared Components stellt Softwarekomponenten bereit, die von den anderen Teilprojekten verwendet werden können.
\item Ant Tasks ermöglicht es, Maven-Funktionalität aus Ant-Skripten heraus zu verwenden.
\item Doxia ist ein Framework zum Generieren von Content aus den Formaten Almost Plain Text (APT), Confluence, DocBook, FML (FAQ Markup Language), LaTeX, Markdown, Rich Text Format (RTF), TWiki, XDoc und XHTML.
\item SCM (Source Code Management) entwickelt Software für die Anbindung von Apache an verschiedene Systeme zur Versionsverwaltung wie CVS oder Subversion.
\item Surefire entwickelt ein Testframework für Maven.
\item Wagon stellt eine Abstraktionsschicht für Kommunikationsprotokolle wie "Dateizugriff“, HTTP oder FTP bereit.
\end{itemize}

\subsection{Versionsgeschichte}
Die erste Version, Maven 1.x, wurde im Jahre 2003 eingeführt und am 13. Juli 2004 als Version 1.0 fertiggestellt. Die Umsetzung passierte jedoch sehr schnell, sodass einige Eigenheiten nicht bedacht wurden. Beispielsweise gibt es Performanceprobleme und zu viele Konfigurationsdateien und -angaben.

Deshalb wurde das Konzept überarbeitet und seit dem Jahre 2005 parallel begonnen, Maven 2.x zu entwickeln, welches in Version 2.0 am 19. Oktober 2005 fertiggestellt wurde.

Maven 1.x wird nicht mehr weiterentwickelt und beschränkt sich auf Support und Bugfixes. Die letzte Version von Maven 1.x (Version 1.1) wurde am 25. Juni 2007 ausgeliefert.

Die Entwicklung von Maven 3.0 begann im Jahr 2008. Nach acht Alpha-Releases wurde die erste Beta-Version von Maven 3.0 im April 2010 veröffentlicht. Besonderes Augenmerk lag auf der Kompatibilität zwischen Maven 2 und 3.

